% Pista alg-26s-q2 · Àlgebra
% Procedència: PAU Catalunya 2026 · Convocatòria de setembre · Sèrie 2 · Exercici 2
\documentclass[11pt,a4paper]{article}
\usepackage{pista}
\pistainfo{Catalunya 2026}{Convocatòria de setembre}{Sèrie 2}

\begin{document}

\section*{\textcolor{blauPista}{Exercici 2}}

\noindent Considereu el sistema d'equacions següent:
\[
\left.\begin{array}{r}
2x + ky + z = 3 \\
-x + 2y + 2z = 1 \\
x + y + z = 2
\end{array}\right\},
\]
on $k$ és un paràmetre real.

\subsection*{2.1. \normalfont Discutiu el sistema en funció del paràmetre $k$. \hfill\textnormal{[1~punt]}}

\pista{Escriu la matriu de coeficients $A$ i la matriu ampliada $A^{*}$ i calcula $\det(A)$ (vigila els signes dels termes de Sarrus que porten $k$): t'ha de sortir una expressió de primer grau en $k$, que s'anul·la per a un sol valor. Per a tots els altres valors, aplica el teorema de Rouché-Frobenius. Per al valor crític, no donis per fet que el sistema és incompatible: busca un menor d'ordre 2 no nul de $A$ i orla'l amb la fila que hi falta i la columna dels termes independents per decidir el rang de $A^{*}$.}

\subsection*{2.2. \normalfont Resoleu el sistema per a tots els valors de $k$ i comproveu que, per a $k \neq 1$, la solució no depèn de $k$. \hfill\textnormal{[1~punt]}}

\pista{Hi ha dos casos. Per a $k \neq 1$, aplica Cramer amb $\det(A)$ com a denominador i factoritza els numeradors: pots simplificar el factor que depèn de $k$ precisament perquè $k \neq 1$. La tercera incògnita, treu-la de l'equació més senzilla. Substituint la solució a la primera equació veuràs per què $k$ no hi influeix. No oblidis el cas $k = 1$ (et demanen tots els valors): queda't amb dues equacions independents, pren una incògnita com a paràmetre $\lambda$ i expressa-hi les altres dues.}

\subsection*{2.3. \normalfont Interpreteu geomètricament el resultat dels apartats anteriors. \hfill\textnormal{[0,5~punts]}}

\pista{Cada equació és un pla de l'espai. Fixa't que la segona i la tercera no depenen de $k$: són dos plans fixos que es tallen en una recta $r$, i només el primer pla canvia amb $k$. Tradueix la discussió a posicions relatives: què vol dir, per a tres plans, un sistema compatible determinat? I un de compatible indeterminat amb un grau de llibertat (comprova que no hi ha dos plans coincidents)? Quina relació hi ha, doncs, entre el primer pla i $r$ quan $k = 1$? Finalment, interpreta que el punt solució no depengui de $k$: per on passen tots els plans de la primera equació? Comprova que aquest punt també és sobre $r$.}

\end{document}
